\section{Introduction}
\label{sec:introduction}
After a multi-year, multi-round public competition, NIST published its first post-quantum standards FIPS-203~\cite{fips203}, FIPS-204~\cite{fips204}, 
and FIPS-205~\cite{fips205} in 2024, specifying ML-KEM, ML-DSA, and SLH-DSA, respectively. Later in 2025, NIST selected an additional KEM, HQC, 
for standardization in the fourth round. Unlike ML-KEM, which is based on module-LWE, HQC relies on the Quasi-Cyclic Syndrome Decoding (QCSD) 
problem ~\cite{hqc2025spec}, providing 
diversity beyond lattice-based assumptions. This choice also marked a shift away from the lattice-based KEMs that dominated the 
Round 3 submissions. Four of the schemes in the PQC Round
3 Submissions \cite{nistpqc-round3} -- \textsc{Kyber} \cite{kyberoriginal} (standardized as ML-KEM), Saber \cite{AFRICACRYPT:DKRV18}, NTRU LPRime \cite{ntruprime}, and FrodoKEM \cite{takeofffrodo} --
are closely related lattice KEMs, spanning a
spectrum of algebraic structure, from FrodoKEM's unstructured LWE, through \textsc{Kyber}'s and Saber's cyclotomic modules, to
the prime-degree number fields of NTRU Prime. This spectrum reflects different design principles, commonly understood as
structure vs. conservatism. ML-KEM's module structure gives small keys and fast NTT arithmetic, 
but is also seen as a potential attack surface, so conservative designs discard it entirely \cite{takeofffrodo}. FrodoKEM, for example,
is based on a more conservative hardness assumption of plain LWE, which is why bodies such as BSI, ENISA, and NLNCSA point
to it for long-term confidentiality \cite{frodokem-news}; FrodoKEM was also recently standardized by ISO \cite{frodoiso}.

It is useful to separate the two notions of conservatism this framing bundles together as they concern two different
(and today, entirely hypothetical) classes of attack. The first arises from the choice of hardness assumption. 
FrodoKEM rests on unstructured LWE to avoid a hypothetical structure break, an attack that would exploit the algebraic structure of 
ring/module LWE but not apply directly to plain LWE.
The second involves the security margin a scheme carries above its target level, which protects against advances 
in the best known attacks today. This could include improved lattice reduction, dual attacks, or sharper cost models 
that would affect the concrete security of every lattice scheme at its current parameters, including FrodoKEM. Both of these 
hedges are distinct, but not entirely independent: additional parameter margin can also increase the cost of a 
hypothetical structural attack, while removing structure can affect the cost of known attacks. A third, orthogonal hedge is to change the hardness 
assumption altogether, as HQC does by relying on QCSD instead of LWE. This provides cryptographic diversity, since an attack on one assumption 
does not automatically imply an attack on the other. 
However, diversity alone also does not imply that the schemes are cryptanalytically independent.
Recent work has demonstrated connections between lattice and code-based cryptanalysis, including adaptations of lattice-sieving techniques to codes \cite{EC:DEEK24}.
Moreover, the quasi-cyclic structure of HQC has been noted to have important parallels with the ring and module structures used in 
structured lattice schemes \cite{hqc-forum}. 
These approaches address different sources of unknowns, and it is impossible to predict which is more likely.

Faced with this uncertainty, increasing the margin has been a historically standard response, provided that sufficient bandwidth is available.
For computationally secure schemes,
a concrete security estimate is the cost of the best attack known at the time, and such estimates evolve as cryptanalysis matures.
For instance, the recommended size of an RSA modulus climbed from $512$ to $1024$ to $2048$ bits and beyond as factoring 
algorithms and available computation advanced \cite{nist-sp800-57pt1r5}. Lattice estimates are not static either. The cost models used for
estimation and the attacks have evolved over more than a decade \cite{bkz2,estimator}. The core framework for estimating 
the security of LWE-based schemes has nevertheless remained stable, with BKZ-based lattice reduction at its core. Within this framework, the cost of the best known 
attacks continues to drift in both directions as improvements are proposed \cite{fasterdual,matzov2022lwe}, scrutinized 
\cite{doesitwork}, and refined \cite{variantofmatzov}.

If the best known attacks continue to improve, how much margin is required remains open, and the useful question is not whether 
today's KEMs carry enough but how much margin they can carry at all. This reframing matters because additional margin comes at a cost, 
requiring larger parameters and hence higher bandwidth and runtime, so no scheme can indefinitely withstand advances in 
cryptanalysis simply by scaling its parameters. For LWE-family KEMs there is a further constraint, since the parameters 
that determine security also control the decryption failure probability, so increasing the noise to gain a larger security 
margin necessarily affects correctness.

Crucially, what structure determines is the \emph{price} of buying additional margin. At matched security levels, an unstructured scheme 
already requires roughly an order of magnitude more bandwidth than a module scheme \cite{takeofffrodo,kyberoriginal}. Since these communication costs are 
considered practical enough for conservative deployments, a module scheme can spend the same bandwidth budget on larger parameters and obtain
a conservative security margin while retaining the efficiency advantages of NTT arithmetic. 

Recent work has explored this high-margin regime from the perspective of unstructured LWE. KyFrog \cite{kyfrog} increases the LWE dimension and 
adjusts the noise distribution to achieve substantially higher security estimates for much larger ciphertexts. This motivates the complementary question of 
how much security margin a module KEM can achieve within the communication budget of a conservative unstructured LWE scheme, and at what computational cost. 
We answer this question with \textsc{Kaiburr}, a module KEM that scales ML-KEM's 
module rank and introduces a generalized noise distribution to preserve correctness. We first characterize 
the correctness and security limits of naively increasing ML-KEM's module rank, then introduce a noise family that 
allows further scaling. We then build three different parameter sets with the new noise sampler, and implement them across conventional software and
dedicated hardware-accelerated platforms. Finally, we compare \textsc{Kaiburr} with existing KEMs, including FrodoKEM 
and HQC. More specifically, our contributions can be summarized as follows.

\subsection{Our contributions}
\begin{itemize}
  \item \emph{We show that naively scaling ML-KEM's module rank hits a correctness wall.} Raising $k$ alone, with \textsc{Kyber}'s noise and 
    modulus fixed, satisfies the $2^{-128}$ decryption-failure
    bound only up to $k=7$ and fails already at $k=8$, so larger module rank cannot buy additional
    security margin on its own (Section~\ref{sec:warmup}).
  \item \emph{We introduce a one-parameter noise family that allows us to scale further.} Our family $f(t)$
    generalizes \textsc{Kyber}'s centered binomial distribution, thinning its $\pm 2$ tail geometrically as $t$
    grows while leaving the $\pm 1$ mass unchanged.
    The parameter $t$ can be tuned against module rank $k$ to keep decryption failure below
    $2^{-128}$ as $k$ increases (Section~\ref{sec:distribution}).
  \item \emph{We select concrete parameter sets that reach large security margins within FrodoKEM-640's
    communication budget.} Using this trade-off, we obtain three parameter sets, kaiburr-1792, -4608, and
    -6144, at $k=7,18,24$, whose classical/quantum Core-SVP security climbs from $479/434$ to $1717/1557$
    bits while their public keys and ciphertexts stay within the communication budget of FrodoKEM-640, its smallest standardized
    parameter set (Section~\ref{sec:benchmarks}).
  \item \emph{We build \textsc{Kaiburr} across three implememtation targets.} We implement all three
    parameter sets in C, in Jasmin with machine-checked constant-time verification,
    and as an ACC variant targeting the Pavona hardware root of trust, detailing the engineering
    consequences of scaling $k$ and removing ciphertext compression (Section~\ref{sec:implementation}).
  \item \emph{We give a formally verified bound on Kaiburr's decryption-failure probability.} 
  Building on the EasyCrypt framework of~\cite{BarbosaKLSS25}, we show that removing ciphertext compression is what 
  makes this bound provable rather than heuristic and close a gap that exists even for ML-KEM's own published bound 
  (Section~\ref{sec:verification}).
  \item \emph{We benchmark \textsc{Kaiburr} against FrodoKEM and HQC.} Across three Intel
    microarchitectures, we show that our AVX2 and Jasmin implementations run substantially faster than FrodoKEM-640 and 
    achieve performance comparable to HQC, despite having (much) larger security margins. (Section~\ref{sec:benchmarks}).
\end{itemize}
\subsection{Access to artifacts}
All software, including parameter analysis scripts, implementation source code, and other artifacts, is publicly available on \href{https://github.com/cryptojedi/kaiburr}{GitHub}.